% !PRINTED 240
\weberpage{240}
%<<CONT
$j(\omega)$ mit beliebigen Zahlenkoeffizienten betrachtet wird. Denn
besteht irgend eine rationale Gleichung
\weq{15}{\Phi[j(n\omega),\ j(\omega)] = 0,}
so darf darauf jede beliebige lineare Transformation
\[
\begin{pmatrix} \alpha, & \beta \\ \gamma, & \delta \end{pmatrix}
\]
angewandt werden, und nach § 29, (5) lassen sich, wenn
\[
\begin{pmatrix} a, & b \\ c, & \partial \end{pmatrix}
\]
eine beliebige Transformation $n$ter Ordnung ist, die linearen
Transformationen
\[
\begin{pmatrix} \alpha, & \beta \\ \gamma, & \delta \end{pmatrix}, \quad
\begin{pmatrix} \alpha', & \beta' \\ \gamma', & \delta' \end{pmatrix}
\]
immer so bestimmen, daß
\weq{16}{\begin{pmatrix} \alpha', & \beta' \\ \gamma', & \delta' \end{pmatrix}\begin{pmatrix} a, & b \\ c, & \partial \end{pmatrix} = \begin{pmatrix} 1, & 0 \\ 0, & n \end{pmatrix}\begin{pmatrix} \alpha, & \beta \\ \gamma, & \delta \end{pmatrix}.}
% [sic] The source says equation (9), but the argument clearly refers to equation (15).
Daraus folgt aber, daß die Gleichung (9) durch \gesperrt{jede} der
Größen
\[
j\left(\frac{c+\partial\omega}{a+b\omega}\right)
\]
und mithin auch durch jede der Größen (1) befriedigt ist, woraus
[wegen der Verschiedenheit der Größen (1)] die Irreducibilität
von (14) folgt.

2. Die Funktion
\weq{17}{F_n[x, j(\omega)] = \Pi\left[x - j\left(\frac{c+\partial\omega}{a}\right)\right]}
hat für jeden endlichen Wert von $\omega$ mit positiv imaginärem Bestandteil, also auch für jeden \gesperrt{endlichen Wert} von $j(\omega)$, einen endlichen Wert und ist sonach eine \gesperrt{ganze} rationale Funktion von $j(\omega)$.

Suchen wir ferner nach (3) das Anfangsglied der Entwickelung
der Funktion (17) nach Potenzen von $q$, so ergibt sich
\[
(-1)^{\nu} e^{-2\pi i \Sigma \frac{c}{a}} e^{-2\pi i \omega \Sigma \frac{\partial}{e}\varphi(e)} = C q^{-2\nu},
\]
wenn $C$ eine endliche Konstante ist. Es ist daher
\[
j(\omega)^{-\nu} F_n[x, j(\omega)]
\]
für ein unendliches $j(\omega)$ endlich, d. h. der Grad von $F_n[x, j(\omega)]$
in bezug auf $j(\omega)$ ist ebenfalls der $\nu$te.
